\section{Experiments}
\label{sec:experiments}
\subsection{Experimental Setup}
\label{subsec:exp_setup}
\textbf{Models.  }We evaluate a diverse set of general-purpose multimodal large language models (MLLMs) that can (1) process input resolution greater or equal to $448 \times 448$ and (2) achieve a score of at least 36 on the \textit{testmini} set of MathVista \cite{lu2024mathvista}.
For open-source models, we test: InternVL Chat V1.5 \cite{chen2024far}, InternLM-XComposer2-4KHD (IXC2 4KHD) \cite{dong2024internlmxcomposer24khd}, InternLM-XComposer2 (IXC2) \cite{chen2023internvl}, LLaVA 1.6 Yi 34B \cite{liu2024llavanext}, LLaVA 1.6 Mistral 7B \cite{liu2024llavanext}, DeepSeek VL \cite{lu2024deepseekvl}, MoAI \cite{lee2024moai}, IDEFICS 2 \cite{laurenccon2024matters}, IDEFICS 2 Chatty \cite{laurenccon2024matters}, SPHINX V2 \cite{gao2024sphinxx}, Mini-Gemini (MGM) HD Yi 34B \cite{li2024minigemini}, Mini-Gemini (MGM) HD LLaMA3 8B \cite{li2024minigemini}, and MiniCPM-V2 \cite{viscpm} (See more model details in \cref{tab:os_model_components}). 
We also evaluate the following proprietary models: GPT-4o \cite{openai2024gpt4}, GPT-4V \cite{openai2024gpt4}, Claude-3 Opus \cite{claude3}, Claude 3 Sonnet \cite{claude3}, Claude 3 Haiku \cite{claude3}, Reka Core \cite{rekateam2024reka}, Reka Flash \cite{rekateam2024reka}, Reka Edge \cite{rekateam2024reka}, Gemini 1.0 Pro \cite{team2023gemini}, Qwen VL Plus \cite{bai2023qwenvl}, and Qwen VL Max \cite{bai2023qwenvl}.
For all models, we provide generation configurations in \cref{tab:run_configurations}.

\textbf{Baselines.  }We provide a text-only baseline, denoted as Random (GPT-4o), where we prompt GPT-4o to reasonably guess the answer without seeing the charts (see the prompt in \cref{sec:random_guess_prompt}). 
This accounts for the effect of using common sense or shallow cues in textual queries to correctly guess the answer.
We also recruit in-house human participants and report their performance (\textit{i.e.,} Human) on \ours{}.
Notably, we ensure that the participants see the exact same questions and instructions as the models and that their responses are evaluated in the same way as the models' responses. This approach allows us to fairly compare the performance gap between humans and models.

\subsection{Experimental Results}
\label{subsec:exp_results}
We provide quantitative results on the validation set for all models in \cref{tab:main_result}. Additional results on the test set are available in \cref{tab:test_set_full_result}. To better understand where models fail, we select a set of representative models \cite{openai2024gpt4, claude3, rekateam2024reka, chen2024far, li2024minigemini, laurenccon2024matters} and present examples of failure cases for 30 descriptive questions in \cref{sec:desc_failure} and 30 reasoning questions in \cref{sec:reas:failure}. The latest results are in our \href{https://charxiv.github.io/#leaderboard}{leaderboard}.

\begin{table*}
  \centering
  \caption{Evaluation results on the validation set. Bold numbers represent the best in-class performance (open-source or proprietary), and \underline{underlined} numbers represent the second-place. Models with (*) are those whose performance is constrained by input resolutions (see \cref{tab:os_model_components} for details). Info. Extr.$=$information extraction, Enum.$=$enumeration, Patt. Rec.$=$pattern recognition, Cntg.$=$counting, Comp.$=$compositionality. Details for these categories are shown in \cref{fig:overview,,subsec:charxiv_questions}.}
  \scalebox{0.68}{
  \begin{tabular}{llarrrrracccccc@{}}
    \hline
    \toprule
    \multirow{4}{*}{\textbf{Model}} & \phantom{} & \multicolumn{5}{c}{\textbf{Reasoning Questions}} & \phantom{} & \multicolumn{6}{c}{\textbf{Descriptive Questions}} & \phantom{} \\
    

    \cmidrule{3-7} \cmidrule{9-14} 

    \phantom{} & \phantom{} &  & \small \textbf{Text in} & \small \textbf{Text in} & \small \textbf{Num. in} & \small \textbf{Num. in}  & \phantom{} &  & \small \textbf{Info.} & \small \textbf{Enum.} & \small \textbf{Patt.} & \small \textbf{Cntg.} & \small \textbf{Comp.}   \\

    \phantom{} & \phantom{} & \multirow{-2}{*}{\textbf{All}} & \small \textbf{Chart} &  \small \textbf{General} & \small \textbf{Chart} & \small \textbf{General}  & \phantom{} & \multirow{-2}{*}{\textbf{All}} & \small \textbf{Extr.} & \small \textbf{} & \small \textbf{Rec.} & \small \textbf{} & \small \textbf{}   \\
    \midrule
            \multicolumn{14}{c}{\textbf{Baselines}} \\
    \midrule
Human   &  & \textbf{80.50} & \textbf{77.27} & \textbf{77.78} & \textbf{84.91} & \textbf{83.41} &  &  \textbf{92.10}               &   \textbf{91.40}            &   \textbf{91.20}            &   \textbf{95.63}          &  \textbf{93.38}           &   \textbf{92.86}            \\
Random (GPT-4o) \cite{openai2024gpt4}         &  & 10.80 &  4.32 & 39.39 & 5.60 & 16.16 &  &   19.85             &    21.65           &    16.71             &  23.80             &      25.70          &   5.36              \\
    \midrule
        \multicolumn{14}{c}{\textbf{Proprietary Multimodal Large Language Models}} \\
    \midrule
GPT-4o \cite{openai2024gpt4}         &  & \textbf{47.10} & \textbf{50.00} & \textbf{61.62} & \textbf{47.84} & \textbf{34.50} &  & \textbf{84.45}                & \textbf{82.44}                & \textbf{89.18}                & \textbf{90.17}                & \textbf{85.50}                &\textbf{59.82}                \\
GPT-4V  \cite{openai2024gpt4}        &  & \underline{37.10} & \underline{38.18} & \underline{57.58} & \underline{37.93} & 25.33 &  & \underline{79.92}                & \underline{78.29}                & \underline{85.79}                & \underline{88.21}                & \underline{80.92}                & \underline{41.07}                \\
Claude 3 Sonnet \cite{claude3} &  & 32.20 & 31.59 & 50.51 & 31.47 & 26.20 &  & 73.65                & 75.74                & 81.92                & 76.64                & 72.26                & 8.48                 \\
Claude 3 Haiku \cite{claude3} &  & 31.80 & 29.77 & 45.45 & 34.48 & \underline{27.07} &  & 65.08                & 69.87                & 69.98                & 64.85                & 61.83                & 8.04                 \\
Claude 3 Opus \cite{claude3}  &  & 30.20 & 26.36 & 50.51 & 33.62 & 25.33 &  & 71.55                & 75.62                & 73.69                & 73.58                & 70.48                & 26.79                \\
Reka Core \cite{rekateam2024reka}      &  & 28.90 & 27.50 & 41.41 & 28.45 & 26.64 &  & 55.60 &	58.90&	50.52&	65.72&	71.25&	10.71\\
Reka Flash \cite{rekateam2024reka}     &  & 26.60 & 26.59 & 39.39 & 30.60 & 17.03 &  & 56.45                & 61.39                & 48.59                & 69.87                & 72.52                & 7.14                 \\
Qwen VL Max \cite{bai2023qwenvl}    &  & 24.70 & 26.14 & 41.41 & 24.57 & 14.85 &  & 41.48                & 50.42                & 28.41                & 53.71                & 51.15                & 4.46                 \\
Reka Edge \cite{rekateam2024reka}      &  & 23.50 & 20.23 & 32.32 & 30.60 & 18.78 &  & 33.65                & 36.65                & 28.49                & 34.72                & 52.16                & 4.91                 \\
Gemini 1.0 Pro \cite{team2023gemini}  &  & 22.80 & 20.91 & 48.48 & 18.10 & 20.09 &  & 54.37                & 67.97                & 39.23                & 60.48                & 62.60                & 8.93                 \\
Qwen VL Plus \cite{bai2023qwenvl}    &  & 16.00 & 15.45 & 45.45 & 12.07 & 8.30  &  & 28.93                & 33.33                & 17.92                & 32.10                & 56.23                & 2.23                 \\
    \midrule
        \multicolumn{14}{c}{\textbf{Open-Source Multimodal Large Language Models}} \\
    \midrule
InternVL Chat V1.5 \cite{chen2024far}       &  & \textbf{29.20} & \textbf{30.00} & \textbf{45.45} & \textbf{32.33} & 17.47 &  & \textbf{58.50} & \textbf{69.63} & 52.95 & 53.06 & \textbf{64.63} & 5.80 \\
MGM HD Yi 34B \cite{li2024minigemini}    &  & \underline{25.00} & \underline{26.59} & \underline{43.43} & 27.16 & 11.79 &  & 52.68 & 53.86 & \underline{55.04} & \textbf{65.50} & 53.94 & 2.23 \\
IXC2 4KHD \cite{dong2024internlmxcomposer24khd} &  & \underline{25.00} & 23.86 & \underline{43.43} & \underline{29.31} & 14.85 &  & \underline{54.65} & \underline{61.09} & 54.08 & 51.53 & 59.80 & 6.70 \\
LLaVA 1.6 Yi 34B* \cite{liu2024llavanext}        &  & 22.50 & 20.45 & 37.37 & 23.71 & \underline{18.78} &  & 51.05 & 46.38 & \textbf{63.44} & \underline{56.11} & 51.91 & 5.80 &  \\
MGM HD LLaMA3 8B \cite{li2024minigemini} &  & 19.00 & 19.77 & 36.36 & 21.12 & 7.86  &  & 44.42 & 49.41 & 39.23 & 51.09 & 55.98 & 1.79 \\
IXC2* \cite{chen2023internvl}      &  & 18.70 & 16.14 & 38.38 & 21.98 & 11.79 &  & 38.75 & 34.10 & 43.58 & 46.72 & 52.93 & 5.80 \\
MiniCPM-V2 \cite{viscpm}               &  & 18.50 & 17.95 & 33.33 & 19.40 & 12.23 &  & 35.77 & 39.74 & 36.56 & 26.42 & 44.53 & 5.36 \\
IDEFICS 2 \cite{laurenccon2024matters}                &  & 18.20 & 15.45 & 35.35 & 17.24 & 17.03 &  & 32.77 & 36.12 & 27.28 & 40.83 & 43.26 & 3.12 \\
IDEFICS 2 Chatty \cite{laurenccon2024matters}         &  & 17.80 & 15.45 & 34.34 & 19.83 & 13.10 &  & 41.55 & 34.88 & 54.56 & 45.63 & 44.27 & 6.70 \\
MoAI* \cite{lee2024moai}                    &  & 17.50 & 9.32  & 36.36 & 21.12 & \textbf{21.40} &  & 28.70 & 31.20 & 21.23 & 39.96 & 40.46 & \underline{7.59} &  \\
DeepSeek VL \cite{lu2024deepseekvl}             &  & 17.10 & 16.36 & 32.32 & 19.83 & 9.17  &  & 45.80 & 49.11 & 45.20 & 42.79 & \underline{60.31} & 4.91 \\
SPHINX V2* \cite{gao2024sphinxx}               &  & 16.10 & 13.86 & 28.28 & 17.67 & 13.54 &  & 30.25 & 35.59 & 24.37 & 41.05 & 29.52 & 1.79 \\
LLaVA 1.6 Mistral 7B* \cite{liu2024llavanext}   &  & 13.90 & 11.36 & 32.32 & 16.81 & 7.86  &  & 35.40 & 34.70 & 33.98 & 48.91 & 42.49 & \textbf{8.48} \\
    \bottomrule
    \hline
  \end{tabular}
  }
  \vspace{-2ex}
  \label{tab:main_result}
\end{table*}

\textbf{All models struggle with reasoning questions.  } 
As shown in \cref{tab:main_result}, the top-performing model, GPT-4o, only correctly answers $47.1$\% of the reasoning questions, exhibiting a $33.4\%$ gap to the human performance of $80.5\%$. 
Moreover, the strongest open-source model, InternVL Chat V1.5, only correctly answers $29.2\%$ of the reasoning questions, highlighting a substantial gap between the leading proprietary and open-source model.
Notably, none of the other open-source models can correctly answer more than $25\%$ of the reasoning questions, indicating marked weaknesses in handling the diverse and challenging chart reasoning questions in \ours{} despite achieving decent performance in existing benchmarks \cite{kafle2018dvqa, kahou2017figureqa, masry2022chartqa, lu2024mathvista} (\textit{e.g.,} see \cref{fig:benchmark_comparison}).

\textbf{Open-source models still struggle with descriptive questions.  }
The leading proprietary model, GPT-4o, exhibits strong capabilities in answering descriptive questions, lagging just $7.65\%$ behind human performance.
However, similar to our findings on reasoning questions, the top-performing open source model InternVL Chat V1.5 exhibits a $25.95\%$ drop in performance compared to GPT-4o.
Overall, the performance of open-source models on descriptive questions remains very low, with most models failing to correctly answer more than $50$\% of questions.

\subsection{Analysis}
\label{subsec:exp_analysis}

\textbf{Descriptive skills are a prerequisite for reasoning.  }
We find that models with strong reasoning capabilities exhibit strong descriptive capabilities, but the reverse is \emph{not} guaranteed (e.g., see Gemini 1.0 Pro, IDEFICS 2 Chatty and DeepSeek VL in Tab. \ref{tab:main_result}).
Manual inspection of models' answers to reasoning questions reveals that some models \cite{rekateam2024reka, li2024minigemini, bai2023qwenvl, lee2024moai} leverage zero-shot Chain-of-Thought (CoT) reasoning \cite{wei2022chain, zhang2023multimodal} to answer the reasoning questions.
However, such CoT may not always be helpful, especially when models cannot accurately describe the chart, as we show in \cref{subsec:desc_example13,,subsec:desc_example28,,subsec:reas_example1,,subsec:reas_example17}.
%where errors in describing elements of the chart often lead to incorrect answers to the reasoning questions.
Quantitatively, we show in \cref{sec:rel_gen_corr} that more lengthy responses (\textit{e.g.,} those potentially containing more CoT traces) can \textit{negatively} impact models' performance on reasoning questions. 
This issue is especially clear in models with low accuracy on descriptive questions, such as MoAI and Qwen VL Plus, which answer $28.70$\% and $28.93$\% of descriptive questions correctly.
In contrast, models with higher accuracy on descriptive questions, such as Mini-Gemini HD Yi 34B and Reka Flash, which achieve $52.68$\% and $56.45$\%, respectively, show improved performance on reasoning questions when generating lengthy responses.
Nevertheless, the vast majority of models exhibit performance uncorrelated with response length.
Thus, we hypothesize that a model must have a strong basic understanding in order to generate helpful multimodal CoT for reasoning.

\textbf{Models struggle with compositional tasks that are easy for humans.  } 
We find that the descriptive task that most strongly differentiates the capabilities of the leading open-source, the top-performing proprietary model, and humans is to count the number of labeled ticks on the x- and y-axes (see \cref{subsec:desc_example28}), on which they achieve $92.86$\%, $59.82$\% and $5.80$\% accuracy respectively.
Although counting is easy for humans, this particular task causes 20 out of 24 models to achieve an accuracy below $10$\% (our random baseline achieves $5.35$\%).
While we do not specifically measure how close each model's responses are to the ground truth, a near-random performance pinpoints the weakness of MLLMs in solving compositional and novel chart understanding tasks.

\begin{figure}[t!]
\centering
\subfigure[]
{\label{fig:hallucinations}\includegraphics[width=100mm]{figures/descriptive_hallucinations.pdf}}
\hspace{0mm}
\subfigure[]{\label{fig:num_subplots}\includegraphics[width=35mm]{figures/num_subplots.pdf}}
\vspace{-2ex}
\caption{Analysis on unanswerable questions (a) and charts with subplots (b).}
\vspace{-4ex}
\end{figure}

\textbf{Weak models cannot identify unanswerable questions.  }
\ours{} is the first work to introduce unanswerable questions in chart understanding. 
As discussed in \cref{subsec:charxiv_questions}, $25$\% of descriptive questions are designed to be unanswerable, where the requested information does not exist or is not applicable to the target subplot in the chart (see \cref{subsec:desc_example2,,subsec:desc_example4,,subsec:desc_example6,,subsec:desc_example12,,subsec:desc_example14,,subsec:desc_example16,,subsec:desc_example18,,subsec:desc_example20,,subsec:desc_example22,,subsec:desc_example24,,subsec:desc_example26}).
We measure how often models can correctly identify and suitably respond to unanswerable questions in \cref{fig:hallucinations}.
Interestingly, the models that achieve an accuracy below $80$\% on unanswerable questions each exhibit idiosyncratic patterns of failure.
For example, IDEFICS 2 Chatty incorrectly responds to nearly $90$\% unanswerable questions about the title, x- and y-axis labels, yet correctly identifies more than $90$\% of unanswerable questions about intersections of lines and the presence of the legend.
On the other hand, IXC 2 correctly respond to $80$\% questions about names of title, x- and y-axis labels that are unanswerable, yet fails to identify unanswerable cases for the difference in tick values when ticks are categorical or the difference is not constant. 

\textbf{Descriptive capabilities degrade with more subplots.  }
\ours{} is the first work to aggregate detailed statistics on the number of subplots in each chart, so we are able to conduct a fine-grained analysis of how the performance of proprietary models and open-source models changes with the number of subplots in the chart. 
As shown in Figure \ref{fig:num_subplots}, a representative set of open-source and proprietary models struggle to answer descriptive questions about charts with more subplots. 
With 6+ subplots, the deterioration is $30$\%--$50$\% for open-source models and only $10$\%--$30$\% for proprietary models.
This indicates that all MLLMs are weaker in handling descriptive queries for charts with more subplots, and such performance deterioration is exacerbated in open-source models.
We hypothesize that this is because open-source models are instruction-tuned on chart datasets that do not contain subplots, such as DVQA and ChartQA.
On the other hand, there appears to be no clear correlation between reasoning capabilities and the number of subplots.

\textbf{Model performance varies among different subjects.  }
Although the questions in \ours{} are designed to be answerable without domain-specific knowledge, we measure the models' performance on individual subjects (see \cref{fig:overview}). 
All models show consistently weaker descriptive capabilities on physics-related charts and stronger performance on charts containing electrical engineering and systems science, quantitative finance and economic data (see \cref{tab:val_desc_set_subject}).
On the other hand, models exhibit idiosyncratic reasoning capabilities over different subjects, demonstrating no clear pattern (see \cref{tab:val_reas_set_subject}).
Interestingly, the strongest open-source model, InternVL Chat V1.5 matches GPT-4V in correctly answering $39.26$\% of the reasoning questions from charts in the math domain, but it significantly lags behind in other domains, exhibiting gaps greater than $20\%$ in the physics and electrical engineering and systems science domains.
These patterns suggest that (1) charts from certain domains are inherently difficult for models to describe and (2) there exist unique skills that are required to perform complex reasoning over charts from different domains.
